
\documentclass{article}
\usepackage{colortbl}
\usepackage{makecell}
\usepackage{multirow}
\usepackage{supertabular}

\begin{document}

\newcounter{utterance}

\centering \large Interaction Transcript for game `bbh', experiment `web\_of\_lies', episode 222 with reflexion{-}llama{-}3.1{-}8b{-}instruct.
\vspace{24pt}

{ \footnotesize  \setcounter{utterance}{1}
\setlength{\tabcolsep}{0pt}
\begin{supertabular}{c@{$\;$}|p{.15\linewidth}@{}p{.15\linewidth}p{.15\linewidth}p{.15\linewidth}p{.15\linewidth}p{.15\linewidth}}
    \# & \multicolumn{2}{c}{Player} && \multicolumn{2}{c}{Game Master} \\
    \hline

    \theutterance \stepcounter{utterance}  
    & & & \multicolumn{4}{p{0.6\linewidth}}{
        \cellcolor[rgb]{0.9,0.9,0.9}{
            \makecell[{{p{\linewidth}}}]{
                \texttt{\tiny{[P1$\langle$GM]}}
                \texttt{You have attempted this same task before and did not succeed. These are the plans you wrote after those attempts.} \\
\texttt{Trial \#1: New plan: I failed to correctly evaluate the boolean function because I didn't properly track the implications of each statement. Specifically, I didn't consider the fact that if Leda lies, then Lorine tells the truth, but I didn't account for the fact that Leda's statement is actually true, which means Lorine's statement about Leda lying is false. Next time, I will carefully track the implications of each statement and consider all possible scenarios, including the case where a statement is true or false. I will also make sure to clearly identify the truth values of each statement and their implications on the overall evaluation.} \\
\\ 
\texttt{Trial \#2: New plan: I failed to correctly evaluate the boolean function because I didn't properly track the implications of each statement. Specifically, I didn't consider the fact that if Leda lies, then Lorine tells the truth, but I didn't account for the fact that Leda's statement is actually true, which means Lorine's statement about Leda lying is false. However, this time I also failed to consider the fact that if Fidel lies, then Leda's statement about Fidel lying is true, which would imply that Lorine's statement about Leda lying is false. Next time, I will carefully track the implications of each statement and consider all possible scenarios, including the case where a statement is true or false. I will also make sure to clearly identify the truth values of each statement and their implications on the overall evaluation. To do this, I will create a truth table to systematically evaluate all possible combinations of truth values for each statement. This will help me to avoid missing any important implications and ensure that I correctly evaluate the boolean function.} \\
\\ 
\\ 
\texttt{Evaluate a random boolean function expressed as a word problem.} \\
\\ 
\texttt{Q: Question: Fidel tells the truth. Jerry says Fidel tells the truth. Vina says Jerry tells the truth. Millicent says Vina lies. Raymond says Millicent lies. Does Raymond tell the truth?} \\
\texttt{A: Yes} \\
\\ 
\texttt{Q: Question: Kristian lies. Millie says Kristian lies. Maybelle says Millie tells the truth. Fidel says Maybelle lies. Leda says Fidel lies. Does Leda tell the truth?} \\
\texttt{A: Yes} \\
\\ 
\texttt{Q: Question: Kristian tells the truth. Michaela says Kristian lies. Raymond says Michaela tells the truth. Osvaldo says Raymond tells the truth. Jamey says Osvaldo tells the truth. Does Jamey tell the truth?} \\
\texttt{A: No} \\
\\ 
\texttt{Q: Question: Teressa lies. Michael says Teressa tells the truth. Fidel says Michael tells the truth. Leda says Fidel lies. Lorine says Leda lies. Does Lorine tell the truth?} \\
\texttt{A:} \\
            }
        }
    }
     \\ \\

    \theutterance \stepcounter{utterance}  
    & \multicolumn{4}{p{0.6\linewidth}}{
        \cellcolor[rgb]{0.561,0.737,0.561}{
            \makecell[{{p{\linewidth}}}]{
                \texttt{\tiny{[P1$\rangle$GM]}}
                \texttt{To evaluate the boolean function expressed in the word problem, I will create a truth table to systematically evaluate all possible combinations of truth values for each statement.} \\
\\ 
\texttt{Let's denote the statements as follows:} \\
\texttt{{-} F: Fidel tells the truth} \\
\texttt{{-} J: Jerry says Fidel tells the truth} \\
\texttt{{-} V: Vina says Jerry tells the truth} \\
\texttt{{-} M: Millicent says Vina lies} \\
\texttt{{-} R: Raymond says Millicent lies} \\
\\ 
\texttt{We are asked to determine if Raymond tells the truth, which is denoted by R.} \\
\\ 
\texttt{First, let's analyze the statements:} \\
\texttt{{-} If Fidel tells the truth (F), then Jerry says Fidel tells the truth (J).} \\
\texttt{{-} If Jerry says Fidel tells the truth (J), then Vina says Jerry tells the truth (V).} \\
\texttt{{-} If Vina says Jerry tells the truth (V), then Millicent says Vina lies (M).} \\
\texttt{{-} If Millicent says Vina lies (M), then Raymond says Millicent lies (R).} \\
\\ 
\texttt{Now, let's create a truth table to evaluate all possible combinations of truth values for each statement.} \\
\\ 
\texttt{| F | J | V | M | R |} \\
\texttt{| {-}{-}{-} | {-}{-}{-} | {-}{-}{-} | {-}{-}{-} | {-}{-}{-} |} \\
\texttt{| T | T | T | T | T |} \\
\texttt{| T | T | T | F | F |} \\
\texttt{| T | T | F | T | F |} \\
\texttt{| T | T | F | F | T |} \\
\texttt{| T | F | T | T | T |} \\
\texttt{| T | F | T | F | F |} \\
\texttt{| T | F | F | T | F |} \\
\texttt{| T | F | F | F | T |} \\
\texttt{| F | T | T | T | F |} \\
\texttt{| F | T | T | F | T |} \\
\texttt{| F | T | F | T | T |} \\
\texttt{| F | T | F | F | F |} \\
\texttt{| F | F | T | T | T |} \\
\texttt{| F | F | T | F | F |} \\
\texttt{| F | F | F | T | F |} \\
\texttt{| F | F | F | F | T |} \\
\\ 
\texttt{From the truth table, we can see that if Fidel tells the truth (F), then Raymond tells the truth (R). However, if Fidel lies (F), then Raymond tells the truth (R) only when Jerry says Fidel tells the truth (J) and Vina says Jerry tells the truth (V).} \\
\\ 
\texttt{Now, let's analyze the statements:} \\
\texttt{{-} If Fidel lies (F), then Jerry says Fidel lies (J).} \\
\texttt{{-} If Jerry says Fidel lies (J), then Vina says Jerry lies (V).} \\
\texttt{{-} If Vina says Jerry lies (V), then Millicent says Vina tells the truth (M).} \\
\texttt{{-} If Millicent says Vina tells the truth (M), then Raymond says Millicent lies (R).} \\
\\ 
\texttt{Now, let's create a truth table to evaluate all possible combinations of truth values for each statement.} \\
            }
        }
    }
    & & \\ \\

    \theutterance \stepcounter{utterance}  
    & & & \multicolumn{2}{p{0.3\linewidth}}{
        \cellcolor[rgb]{0.9,0.9,0.9}{
            \makecell[{{p{\linewidth}}}]{
                \texttt{\tiny{[GM$|$GM]}}
                \texttt{Target: No} \\
            }
        }
    }
    & & \\ \\

    \theutterance \stepcounter{utterance}  
    & & & \multicolumn{2}{p{0.3\linewidth}}{
        \cellcolor[rgb]{0.9,0.9,0.9}{
            \makecell[{{p{\linewidth}}}]{
                \texttt{\tiny{[GM$|$GM]}}
                \texttt{game\_result = LOSE} \\
            }
        }
    }
    & & \\ \\

    \theutterance \stepcounter{utterance}  
    & & \multicolumn{4}{p{0.6\linewidth}}{
        \cellcolor[rgb]{0.95,0.95,0.95}{
            \makecell[{{p{\linewidth}}}]{
                \texttt{\tiny{[P1$\leftrightarrow$P]}}
                \texttt{Post{-}episode Reflexion (not sent to GM):} \\
\texttt{New plan: I failed to correctly evaluate the boolean function because I didn't properly track the implications of each statement. Specifically, I didn't consider the fact that if Fidel lies, then Jerry says Fidel lies, which would imply that Vina says Jerry lies, and then Millicent says Vina tells the truth, and finally Raymond says Millicent lies. However, I also failed to consider the fact that if Fidel lies, then Jerry says Fidel tells the truth, which would imply that Vina says Jerry tells the truth, and then Millicent says Vina lies, and finally Raymond says Millicent lies. Next time, I will carefully track the implications of each statement and consider all possible scenarios, including the case where a statement is true or false. I will also make sure to clearly identify the truth values of each statement and their implications on the overall evaluation. To do this, I will create a truth table to systematically evaluate all possible combinations of truth values for each statement. This will help me to avoid missing any important implications and ensure that I correctly evaluate the boolean function. Additionally, I will make sure to analyze the statements in a more systematic and methodical way, considering the implications of each statement on the overall evaluation, rather than trying to analyze them in a piecemeal fashion.} \\
            }
        }
    }
    & \\ \\

\end{supertabular}
}

\end{document}
